\section{Introduction}
\label{sec:intro}

A decoder-only language model emits one token per forward pass, and pays the
full depth of the network for each one. This is not obviously necessary. Some
continuations are all but fixed by the time the model has read a few tokens:
after \textit{the Empire}, in a sentence about New York, very little other than
\textit{State Building} can follow. If the model's hidden state at
\textit{Empire} already encodes \textit{State Building}, then the two forward
passes that produce those tokens are, in an information-theoretic sense,
re-deriving something the model had already computed.

Recent interpretability work has shown that language models do build
representations of units larger than a single token. \citet{kaplan-etal-2025-tokens}
demonstrate an \emph{intrinsic detokenization} process: when a word is split
into several subword tokens, information from the earlier pieces is aggregated
into the last one, and the resulting state behaves like a representation of the
whole word, even for words the tokenizer never emits as a unit. That result,
however, concerns a unit the model has \emph{already finished reading}. It
establishes that the model consolidates evidence backwards, not that it
anticipates forwards. The question we ask here is the forward one, and it is
the one with consequences for how models are run:

\begin{quote}
Given a partial phrase, does the model's internal representation already encode
the rest, before it has read or generated any of it?
\end{quote}

We study this for \emph{fixed multi-word expressions} --- idioms, landmark
names, and film titles --- because for these the continuation is well defined,
and because their frequency in pretraining data makes them the most favourable
case. If anticipation does not happen here, it is unlikely to happen anywhere.
We index the question by \emph{lookahead distance} $i$: the number of phrase
tokens still unread at the point where we interrupt the model. At $i=4$ in
\textit{The Lord of the Rings}, the model has read only \textit{The}, and we ask
whether \textit{Lord of the Rings} is recoverable from its state.

Recoverable by what means, however, turns out to matter as much as the answer.
We use two readouts on the same interrupted states. The first is
Patchscopes \citep{ghandeharioun-etal-2024-patchscopes}: we lift the hidden
state at the cut position and inject it into an unrelated prompt, and ask
whether the model verbalises the remaining tokens. The second simply lets the
model continue generating from the true prefix it has actually seen. The first
asks what a vector contains; the second asks what the model does. Prior work
treats a successful Patchscopes readout as evidence that the model ``knows''
something, and we wanted to know how far that licence extends.

Our findings are as follows. Phrase-level anticipation is real and substantial:
in natural context, OLMo-2-7B produces the final two tokens of a familiar phrase
in $84.0\%$ of instances, and the final five in $53.6\%$, against $59.2\%$ and
$17.6\%$ for the same phrases in isolation (\S\ref{sec:lookahead}). It is a
stable property of the phrase rather than of the occasion: per-phrase success at
$i=3$ is far more bimodal than independent per-context variation would produce,
with $31.5\%$ of phrases succeeding in at least $90\%$ of their contexts against
a binomial expectation of $10.1\%$ (\S\ref{sec:stability}). And the signal is
decodable: a small
classifier reading the raw hidden state predicts whether the model will in fact
complete the phrase at $76.8\%$ balanced accuracy and $87.6\%$ precision
(\S\ref{sec:probe}).

The two readouts, however, do not measure the same thing. Run on identical
instances, they agree on only $81.1\%$ of them, and their disagreements run in
both directions --- $1{,}924$ instances recovered only by Patchscopes against
$2{,}751$ recovered only by generation (\S\ref{sec:disagreement}). They order
the three phrase categories differently, and the two probes trained on their
respective labels place the ``best'' layer in different parts of the network.
We take this as a caution rather than a refutation: Patchscopes remains
informative about representational content, but a claim of the form \emph{the
model knows $X$ at layer $\ell$} is not readout-independent, and should say
which instrument produced it.

Concretely, we contribute:

\begin{enumerate}[nosep,leftmargin=*]
\item A lookahead-distance protocol for measuring phrase anticipation, applied
      to $1{,}164$ fixed expressions both in isolation and in $9{,}955$
      naturally occurring FineWeb-Edu contexts, with two independent readouts
      run on identical instances (\S\ref{sec:method}).
\item Evidence that anticipation of familiar phrases occurs several tokens
      ahead, is greatly strengthened by natural context, and is a stable
      property of the phrase (\S\ref{sec:lookahead}--\S\ref{sec:stability}).
\item A direct comparison showing that the two readouts disagree substantially
      and asymmetrically, and that conclusions about category difficulty and
      layer localisation depend on which one is used
      (\S\ref{sec:disagreement}).
\item A hidden-state probe that predicts real completion success well above
      chance from the first transformer layer onward, together with an
      adjudication of its apparent errors showing that roughly four in five are
      valid alternate phrasings rather than failures (\S\ref{sec:probe}).
\end{enumerate}
